\documentclass[11pt,a4paper]{article}

\usepackage[utf8]{inputenc}
\usepackage[T1]{fontenc}
\usepackage[spanish,es-tabla]{babel}
\usepackage[margin=2.5cm]{geometry}
\usepackage{lmodern}
\usepackage{amsmath,amssymb}
\usepackage{enumitem}
\usepackage{xcolor}
\usepackage{hyperref}

\hypersetup{
    colorlinks=true,
    linkcolor=blue!70!black,
    urlcolor=blue
}

\begin{document}

\begin{flushleft}
    \textbf{Alumno:} Christian Pérez Hita \\
    \textbf{Asignatura:} Organización y gestión de proyectos \\
    \textbf{Caso práctico:} Proyecto I+D+i
\end{flushleft}

\vspace{0.8em}
\hrule
\vspace{1.5em}

\section*{1. Título del Proyecto}
\noindent\textbf{COOL-SHIELD: Pinturas Inteligentes de Enfriamiento Pasivo y Alta Reflectancia Mediante Materiales Nanocuánticos}

\vspace{1.2em}
\section*{2. Breve Resumen del Proyecto}
El proyecto tiene por objeto la investigación, síntesis y formulación preindustrial de una nueva gama de pinturas acrílicas inteligentes que incorporan aditivos nanocuánticos experimentales. Mediante efectos de confinamiento y dispersión cuántica, estas partículas dotan al recubrimiento de una doble propiedad termo-óptica: maximizar la reflectancia en el espectro solar diurno e inducir emisión térmica radiativa pasiva en la ventana atmosférica infrarroja. De este modo, se logra un enfriamiento sub-ambiente sin consumo de energía ni emisiones directas de CO\textsubscript{2}, optimizando la eficiencia térmica en envolventes arquitectónicas e industriales.

\vspace{1.2em}
\section*{3. Índice de la Estructura del Proyecto de I+D+i}

\begin{itemize}[leftmargin=*, label={}]
    \item \textbf{1. INTRODUCCIÓN Y JUSTIFICACIÓN}
    \begin{itemize}[leftmargin=1.5em, label={}]
        \item 1.1. Contexto industrial y necesidad de recubrimientos atérmicos
        \item 1.2. Principio físico del enfriamiento radiativo pasivo diurno (PDRC)
        \item 1.3. Oportunidad de mercado e innovación tecnológica del producto
    \end{itemize}

    \vspace{0.4em}
    \item \textbf{2. OBJETIVOS DEL PROYECTO}
    \begin{itemize}[leftmargin=1.5em, label={}]
        \item 2.1. Objetivo general
        \item 2.2. Objetivos técnicos y parámetros de formulación
    \end{itemize}

    \vspace{0.4em}
    \item \textbf{3. MATERIALES Y MÉTODOS}
    \begin{itemize}[leftmargin=1.5em, label={}]
        \item 3.1. Selección y funcionalización de nanomateriales cuánticos
        \item 3.2. Formulación de pinturas poliméricas y compatibilización reológica
        \item 3.3. Metodología de caracterización termo-óptica (UV-Vis-NIR y FTIR)
    \end{itemize}

    \vspace{0.4em}
    \item \textbf{4. ACTIVIDADES Y PLAN DE TRABAJO}
    \begin{itemize}[leftmargin=1.5em, label={}]
        \item 4.1. Actividad 1: Síntesis y optimización de nanocomposites
        \item 4.2. Actividad 2: Preparación de lotes piloto y optimización de la dispersión
        \item 4.3. Actividad 3: Ensayos de adherencia, intemperie y envejecimiento acelerado
        \item 4.4. Actividad 4: Validación experimental del enfriamiento sub-ambiente a escala real
    \end{itemize}

    \vspace{0.4em}
    \item \textbf{5. TEMPORALIZACIÓN (CRONOGRAMA)}
    \begin{itemize}[leftmargin=1.5em, label={}]
        \item 5.1. Distribución de actividades y diagrama de Gantt
        \item 5.2. Hitos técnicos y entregables del proyecto
    \end{itemize}

    \vspace{0.4em}
    \item \textbf{6. PRESUPUESTO Y RECURSOS}
    \begin{itemize}[leftmargin=1.5em, label={}]
        \item 6.1. Costes de personal investigador y técnico
        \item 6.2. Costes de materias primas, reactivos y fungibles
        \item 6.3. Amortización de instrumentación y servicios analíticos externos
        \item 6.4. Cuadro resumen del presupuesto de I+D+i
    \end{itemize}

    \vspace{0.4em}
    \item \textbf{7. RESULTADOS ESPERADOS}
    \begin{itemize}[leftmargin=1.5em, label={}]
        \item 7.1. Rendimiento térmico obtenido y reducción de temperatura superficial
        \item 7.2. Viabilidad técnica para el escalado a planta industrial en Kolmer
    \end{itemize}

    \vspace{0.4em}
    \item \textbf{8. PROPIEDAD DE LOS RESULTADOS}
    \begin{itemize}[leftmargin=1.5em, label={}]
        \item 8.1. Estrategia de protección industrial (registro de patente de formulación)
        \item 8.2. Protección del secreto industrial y \textit{know-how} de dispersión
        \item 8.3. Plan de explotación comercial y transferencia tecnológica
    \end{itemize}

    \vspace{0.4em}
    \item \textbf{9. CONCLUSIONES}
\end{itemize}

\end{document}